This thesis examines whether institutional ownership topology improves the prediction of stock-level downside risk beyond conventional market and crowding information. SEC Form 13F holdings and CRSP data are combined into a quarterly stock-level panel covering 2013--2025. Models forecast each stock's maximum drawdown over the following 63 trading days. Ridge and XGBoost test the incremental value of conventional ownership measures and engineered network features, while graph neural networks learn directly from the ownership graph. Models are selected using separate chronological training and validation samples and evaluated on a held-out test period. Predictions from the selected model are converted into a within-quarter Fragility Score, with higher values indicating greater relative downside risk.

A market-only XGBoost model is selected based on its balance of predictive performance and complexity. Its Fragility Score strongly and monotonically ranks realized downside risk throughout the held-out period, although the finalist models perform similarly. Conventional ownership measures provide only small improvements, while engineered network features provide no consistent additional gain. Temporal GraphSAGE performs competitively but does not dominate sufficiently across metrics to justify its additional complexity.

Portfolio tests show that excluding the most fragile stocks modestly improves the downside profile of an equal-weight portfolio, although the result is sensitive to weighting and does not outperform the external market benchmark. An exploratory long--short portfolio further differentiates low- and high-score stocks, but the short test period and incomplete treatment of trading and short-selling costs prevent a tradable-factor interpretation. Overall, the Fragility Score provides an informative forward-looking ranking of equity downside risk, but longer evaluation periods and more realistic implementation assumptions are required before investment deployment.